\documentclass[10pt,a4paper]{amsart}
\usepackage[margin=19mm]{geometry}
\usepackage{amsmath,amssymb,mathtools}
\usepackage[scr=rsfs,cal=euler]{mathalfa}
\usepackage{xfrac}
\usepackage[unicode,hidelinks,bookmarks=false]{hyperref}
\newcommand{\ie}{i.e.\ }
\newcommand{\cf}{cf.\ }
\newcommand{\resp}{respectively\ }
\newcommand{\Subsection}{\subsection}
\providecommand{\nobreakdash}{\nobreak\hbox{-}}
\setlength{\emergencystretch}{2em}
\multlinegap0pt
\usepackage[shortlabels]{enumitem}
\usepackage{tikz}
\usepackage{float}
\usepackage{amssymb}
\usepackage{algorithm}\usepackage{algpseudocode}
\usepackage{xcolor}
\newtheorem{theorem}{Theorem}
\title{Ollivier-Ricci Curvature of Riemannian Manifolds and Directed Graphs with Applications to Graph Neural Networks}
\author{Eleanor P Wiesler}
\date{Source: arXiv:2604.14211v2 (CC BY 4.0)}
\begin{document}
\maketitle

\noindent\emph{Source and licence.} Ollivier-Ricci Curvature of Riemannian Manifolds and Directed Graphs with Applications to Graph Neural Networks. Version arXiv:2604.14211v2 (CC BY 4.0), distributed by arXiv under the Creative Commons Attribution 4.0 International licence, http://creativecommons.org/licenses/by/4.0/, as recorded in the arXiv licence metadata for this exact version and shown on the abstract page https://arxiv.org/abs/2604.14211v2. Full licence notice: \texttt{licenses/arxiv-2604.14211-CC-BY-4.0.txt}.\par\smallskip

\noindent\textit{Editorial scope.} Counted unit: the article's theorem L1-bonnetmyers (source0.tex lines 762--769). The article's complete original proof of it is delivered with it (source0.tex lines 770--817, one \texttt{proof} environment of 48 lines). No mathematics is written, repaired or re-derived: the statement and the proof are the article's own lines, in the article's own notation, and no retained line is substituted or re-flowed.  No further source-local context is needed: every cross-reference of every retained line resolves inside the two delivered spans.  Nothing was omitted from any retained span, so the package carries no omission record.  No retained line contains a citation, so the wrapper carries no bibliography and no bibliography entry.  No retained line contains a figure environment, an image-inclusion command or drawing code, so no image is delivered and the package declares no asset dependency.  Editorial formatting only, in addition: an AMS \texttt{amsart} preamble that restates the article's own declarations and macros used by the retained lines, namely the environments \texttt{theorem} on separate counters, together with the standard packages those lines need.  The retained environments print in this document as Theorem 1 in order of appearance, whereas the article prints the same items with the numbering of its own full document.  The complete native source of the article as distributed in the arXiv e-print for this exact version is bundled unchanged as \texttt{sources/198526/source0.tex}.
\begin{theorem} \label{thm: L1-bonnetmyers}($L^1$ Bonnet-Myers) \\ Suppose that $\operatorname{Ric_O}(x, y) \geq k >0$ for all $x, y \in X$. Then for any $x, y \in X$ we have
$$
d(x, y) \leq  \frac{J(x)+J(y)}{\operatorname{Ric_O}(x, y)}
$$
and therefore
$$\operatorname{diam} X \leq \frac{2 \sup _x J(x)}{k}
$$
\end{theorem}

\begin{proof}
Recall that for a ponint $x \in X$, the expected distance or jump length from x when the random walk begins is defined as follows:
$$
J(x) = \int d(x,y) d \delta_x = W_1(\delta_x, m_x)
$$
Observe that if we compute the 1-Wasserstein distance between two Dirac measures this is equivalent to the distance $d(x,y)$ between the two points these Dirac measures are supported: 
$$
W(\delta_x, \delta_y) = \int d(u,v) d \delta_{(x,y)}(u,v) = d(u,v)
$$
which is intuitive because the cost of transporting a Dirac measure from one point to the other only requires transport from their respective points themselves. From here, we can apply the following triangle inequality as done previously: 
$$
W_1(\delta_x, \delta_y) \leq W_1(\delta_x, m_x) + W_1(m_x, m_y) + W_1(m_y, \delta_y)
$$
which implies that, 
$$
W_1(\delta_x, \delta_y) \leq J(x) + W_1(m_x, m_y) + J(y)
$$
Now recall that we are given the fact that $\operatorname{Ric_O}(x, y) \geq k >0$ for all $x, y \in X$ such that 
$$
1 - \frac{W_1 (m_x, m_y)}{d(x,y)} \geq k >0
$$
So therefore, 
$$
W_1(m_x, m_y) \leq (1-k)d(x,y)
$$
We then obtain the inequality:
\begin{align*}
d(x,y) = W_1(\delta_x, \delta_y &\leq J(x) + W_1(m_x, m_y) + J(y)  \\ &\leq J(x) + J(y) + (1-k)d(x,y)
\end{align*}
and it follows that: 
$$
d(x,y) \leq \frac{\operatorname{Ric}_O}{J(x) + J(x)}
$$
For the second part of this proof we which to show from the previous result that diameter of the space $X$ is bounded by $\frac{2sup_x J(x)}{k}$. By definition: 
$$
\operatorname{diam}X := \operatorname{sup}_{x,y \in X} d(x,y)
$$
Therefore, 
$$
\operatorname{diam} X = \sup_{x, y} d(x, y)
\leq \sup_{x, y} \left( \frac{J(x) + J(y)}{k} \right)
$$
but we arbitrarily taking the supremum for all $x,y$ so the maximal value of $J(x) + J(y)$ will just be $2(J(x))$ for whichever $x$ achieves this upper bound. And we have, 
$$
\operatorname{diam} X \leq 2 \frac{\operatorname{sup}_x J(x)}{k}
$$
as desired. 
\end{proof}

\end{document}
